\scanpage{116}

de $\deleted{\unsure{\Phi}}\Pi_n$ (pol. comb. à $n$ côtés), \deleted{\unsure{au moins}} est
\[
k'_n=\mathbb Z\bigl[\tfrac1n\bigr][U]/\Psi_n[U]
\]
$\alpha'_n\in k'_n$ étant l'image canonique de $U$.

Si $n=\deleted{\unsure{2p}}p$ ($p$ premier\inserted{ impair}\deleted{\unsure{impair}}) \deleted{$p$ et $k$ corps de car. $p$,
alors l'invariant du polygone régulier à $2p$ côtés sur $k$
est $+2$, le polygone cyclotomique}
\deleted{\inserted{D'ailleurs \unsure{si} $p$ impair}}
\[
\deleted{\Phi_{2p}(T)=\Phi_p(-T)\deleted{\unsure{{}^{p}}}
 =\deleted{\unsure{(T^{p-1})(T+1)}}T^{p-1}-T^{p-2}\ldots+1\deleted{\unsure{(T+1)^{p-1}}}}
\]
\deleted{\deleted{\unsure{En car. $p$ on}} En car. $p$, ses solutions sont \deleted{des} toutes
\deleted{$\Psi_{2p}(U)=$} confondues, et égales : $-1$ (celles de $\Phi_p$ étant égales
: $1$) donc les solutions de $\Psi_{\deleted{\unsure{p}}2p}$ sont égales : $1$,
les racines de $\Phi_p(T)=T^{p-1}\ldots1$ \deleted{\unsure{donc alors toutes égales : $1$}}
donc les racines de $\Psi_p(T)$ sont \unsure{toutes} égales : $2$,
Donc si $\alpha\in k$, la représentation de \unsure{donc|dont} celles de
$\Psi_p(U)$ ($=F_p(U)$ ici !) égales : $2$.}
toute réalisation géom. régulière de $\Pi_p$ sur un corps est fidèle
--- d'\deleted{fait}ailleurs on trouve
\[
\begin{aligned}
\Phi_p(T)&=\frac{T^p-1}{T-1}=T^{p-1}\ldots+1\\
\Psi_p(\deleted{\unsure{T}}U)&=F_p(\deleted{\unsure{T}}U)
\end{aligned}
\]
et on trouve l'anneau modulaire
\[
\underline{k}'_p=\mathbb Z[U]/\Psi_p(U)=\mathbb Z[U]/F_p(U)
\]

Supposons enfin que $n=2p$ ($p$ premier). Si $p$
impair, \deleted{\unsure{on a facilement que} ($\zeta\in\overline{\mathbb Q}$) : \unsure{le fait|à la fois}
$(\deleted{\unsure{\ne}}\zeta^2\unsure{-1|\ne1})^p=1$
donc les racines de $\Phi_{2p}(T)=0$ \unsure{sont} les $\zeta$ tels que $(\zeta^2)^p=1$
$\zeta^2\ne1$, $\zeta^p\ne1$ sont aussi les $\zeta\in\overline{\mathbb Q}$ tels que
$(\zeta^2)^p=1$, $\zeta^2\ne1$}
\[
\deleted{\begin{aligned}
 & (T^{2p}-1)\big/\unsure{(T^2-1)(T-1)|(T-1)(T^p-1)}(T^{p-1}\ldots+1)\\
 & \zeta^{2p}=1,\ \zeta^{\unsure{1}}\ne1,\ \zeta^p\ne1,\
 \unsure{(\zeta^2)^p\ne1|\zeta^p\ne1},\ \zeta^2\ne1,\ (T\unsure{+1}),\ \zeta^{2p},\ \zeta^p\ne1
\end{aligned}}
\]
% archive sheet number: 115
